\input{../tex-inputs/latex-leseansicht-vorspann}

\section[Gustav Schwarzkopf an Arthur Schnitzler, 18. 7. 1898]{L04290 Gustav Schwarzkopf an Arthur Schnitzler, 18. 7. 1898}
\nopagebreak\mylabel{L04290v}
\rehead{ }\normalsize\beginnumbering\briefempfaengerindex{Schnitzler, Arthur@\textsc{Schnitzler, Arthur}!zzzSchwarzkopf, Gustav@\emph{von Gustav Schwarzkopf}!1898-07-183@{18. 7. 1898}|(be}
\toendnotes[C]{\smallbreak\pagebreak[2]}
\correspDesc{Versand  durch Gustav Schwarzkopf am 18. 7. 1898 in Wien
\newline{}Übermittlung  am 18. 7. 1898 in Wien
\newline{}Zustellung  am 19. 7. 1898 in Bad Gastein
\newline{}Erhalt  durch Arthur Schnitzler am 19. 7. 1898 in Bad Gastein}\toendnotes[C]{\smallbreak}
\Standort{CUL, Schnitzler, B 96a.}
\physDesc{Postkarte, 913 Zeichen
\newline{}Handschrift: schwarze Tinte, deutsche Kurrent
\newline{}Versand: 1) Stempel: »\nobreak{}\oindex{I., Innere Stadt@\textbf{I., Innere Stadt}, \emph{Verwaltungsgebiet}|pwk}Wien 1/1 1, 18. 7. 98, 8–9N\nobreak{}«.   2) Stempel: »\nobreak{}\oindex{Bad Gastein@\textbf{Bad Gastein}, \emph{Hauptstadt}|pwk}Badgastein, 19/7 \textcolor{gray}{9}8, \textcolor{gray}{5}\nobreak{}«. }\toendnotes[C]{\smallbreak}\pstart{}{\pb}Herrn \textsc{D\textsuperscript{r} Arthur Schnitzler}\pend{}\pstart{}\textsc{Wildbad Gastein\oindex{Bad Gastein@\textbf{Bad Gastein}, \emph{Hauptstadt}|pw}}\pend{}\pstart{}\textsc{Villa Wassing\oindex{Villa Dr. Wassing@\textbf{Villa Dr. Wassing}, \emph{Sanatorium}|pw}}\pend{}{\bigskip}\vspace{1em}
\pstart
           \raggedleft{}{\pb}18. 7. 98\pend
           \vspace{0.5em}
\pstart
           Lieber Arthur, ich werde früheſtens erſt am 26.\textsuperscript{tn.}{ }Abends von Wien\oindex{Wien@\textbf{Wien}, \emph{Verwaltungsgebiet}|pw}
               abreiſen kö{\geminationn}en, weil Max\pwindex{Schwarzkopf, Max 12.\,6.\,1857 Wien – 14.\,4.\,1928 ebd.@\textsc{Schwarzkopf, Max} (12.\,6.\,1857 Wien – 14.\,4.\,1928 ebd.), \emph{Rechtsanwalt}|pw},
               den ich doch gerne erwarten möchte, erſt an dieſem Tage hier eintreffen{ }ſoll. We{\geminationn} Sie am 27.\textsuperscript{ten} von Salzburg\oindex{Salzburg@\textbf{Salzburg}, \emph{Verwaltungsgebiet}|pw} fort wollen, kö{\geminationn}ten wir wol nur einen halben Tag zuſa{\geminationm}en verbringen. Ich ko{\geminationm}e
               aber vielleicht doch, ohne damit ein Verſprechen zu geben. Sie{ }ſind doch in demſelben
                  \textsc{\label{K_L04290-1v}\edtext{Hôtel\oindex{Österreichischer Hof@\textbf{Österreichischer Hof}, \emph{Hotel}|pwv}}{\lemma{\textnormal{\emph{Hôtel}}}\Cendnote{\textnormal{Sie
                     trafen am 4. 8. 1894 in Salzburg\oindex{Salzburg@\textbf{Salzburg}, \emph{Verwaltungsgebiet}|pwk} zusammen. Als Unterkunft
                     dürfte der Österreichische Hof\oindex{Österreichischer Hof@\textbf{Österreichischer Hof}, \emph{Hotel}|pwk} gedient haben, vgl. XXXX Auszeichnungsfehler: Dokument L00361 nicht gefunden.}}}\label{K_L04290-1}}, in dem wir vor vier Jahren gewohnt haben. Selbſtverſtändlich{ }ſind Sie nicht
               gebunden, we{\geminationn} Sie Ihre Dispoſitionen ändern wollen.
               Sollte es{ }ſich nicht. einteilen laſſen, was ich doppelt bedauern würde, weil ich ja
               auch Salzburg\oindex{Salzburg@\textbf{Salzburg}, \emph{Verwaltungsgebiet}|pw} gerne wieder{ }ſehen möchte,{ }ſo
               wünſche ich Ihnen für Ihre Reiſe das Beſte. Auf Wiederſehen da{\geminationn} im Herbſt. – Sie kö{\geminationn}en es
               mir übrigens getroſt glauben, daß es in Wien\oindex{Wien@\textbf{Wien}, \emph{Verwaltungsgebiet}|pw} jetzt{ }ſehr langweilig iſt. – –\pend
           
\pstart
           Mit der Bitte mich Ihrer Frau Mama\pwindex{Schnitzler, Louise 8.\,7.\,1840 Kőszeg – 9.\,9.\,1911 Wien@\textsc{Schnitzler, Louise} (8.\,7.\,1840 Kőszeg – 9.\,9.\,1911 Wien)|pwv}{\\[\baselineskip]}beſtens zu empfehlen –{\\[\baselineskip]}herzlichſt{\\[\baselineskip]}Ihr{\\[\baselineskip]}\spacefill\mbox{Gustav.}\pend
           \leftskip=0em{}\selectlanguage{ngerman}\endnumbering\briefempfaengerindex{Schnitzler, Arthur@\textsc{Schnitzler, Arthur}!zzzSchwarzkopf, Gustav@\emph{von Gustav Schwarzkopf}!1898-07-183@{18. 7. 1898}|)be}\mylabel{L04290h}
\begin{anhang}
\end{anhang}\newcommand{\dateiname}{L04290}\newcommand{\titel}{Gustav Schwarzkopf an Arthur Schnitzler, 18. 7. 1898}\newcommand{\editorInnen}{Herausgegeben von Jahnke, SelmaMüller, Martin Anton}\input{../tex-inputs/latex-leseansicht-abspann}
